% Source-derived chapter 03: Curves and surfaces
% Original source: chabauty-coleman-bound-surfaces
\chapter{Curves and surfaces}
\label{chabauty-coleman-bound-surfaces-chapter-03}
\source{chabauty-coleman-bound-surfaces}

\begin{remark}[Source section]
\label{chabauty-coleman-bound-surfaces-chapter-03-source}
\source{chabauty-coleman-bound-surfaces}
\source{chabauty-coleman-bound-surfaces}
This chapter reproduces the corresponding section of the retrieved
LaTeX source, including its definitions, statements, examples, and
proofs. It has not yet been translated into Lean.
\notready
\end{remark}

% The source section title was: Curves and surfaces
%%%%
%%%%
%%%%
\subsection{Curves} Let $C$ be an irreducible curve over the algebraically closed field $k$ and let $x\in C(k)$. The \emph{order of singularity} of $C$ at $x$ is $\delta(C,x)=\dim_k(\widetilde{\Ocal}_{C,x}/\Ocal_{C,x})$, where the quotient $\widetilde{\Ocal}_{C,x}/\Ocal_{C,x}$ is taken as $k$-vector spaces. One has the following basic fact:
%%
%%
\begin{lemma}
\source{chabauty-coleman-bound-surfaces}
\notready $\delta(C,x)$ is a non-negative integer. It is equal to $0$ if and only if $C$ is regular at $x$.
\end{lemma}
%%
%%

As a special case of Theorem 2 in \cite{Hironaka} we have

%%
%%
\begin{lemma}[Hironaka]
\source{chabauty-coleman-bound-surfaces}
\notready\label{LemmaHironakaGenus} If $C$ is a projective irreducible curve over $k$, then
\source{chabauty-coleman-bound-surfaces}
$$
\gfrak_a(C)=\gfrak_g(C) + \sum_{x\in C(k)} \delta(C,x).
$$

\end{lemma}
%%
%%

The \emph{multiplicity} of $C$ at  $x\in C(k)$ is defined as the multiplicity of the local ring $\Ocal_{C,x}$ as defined, for instance, in Ch.  5, Sec. 14 of \cite{Matsumura}. We denote it by $\mult_x(C)$. A basic fact:
%%
%%
\begin{lemma}
\source{chabauty-coleman-bound-surfaces}
\notready $\mult_x(C)$ is a positive integer. It is equal to $1$ if and only if $C$ is regular at $x$.
\end{lemma}
%%
%%

From Exercise 14.5 in \cite{Matsumura} we get a simple expression for $\mult_x(C)$ in a special case.

%%
%%
\begin{lemma}
\source{chabauty-coleman-bound-surfaces}
\notready\label{LemmaMultFormula} Let $S$ be a surface over $k$ and let $C$ be a closed irreducible curve in $S$. Let $x\in C(k)$ be a point such that $S$ is regular at $x$. Let $F\in \Ocal_{S,x}$ be a local equation for $C$ at $x$. Then 
\source{chabauty-coleman-bound-surfaces}
$$
\mult_x(C)=\max\{n\ge 1 : F\in \mfrak_{S,x}^n\}.
$$
\end{lemma}
%%
%%

We will also need the following result.
%%
%%
\begin{lemma}[Hironaka]
\source{chabauty-coleman-bound-surfaces}
\notready \label{LemmaHironakaMult} Let $S$ be a projective surface and let $C$ be a projective irreducible curve contained in $S$. Let $x\in C(k)$ be a point such that $S$ is regular at $x$. Then we have
\source{chabauty-coleman-bound-surfaces}
$$
\delta(C,x)\ge \frac{1}{2}\cdot \mult_x(C)\cdot (\mult_x(C)-1).
$$
\end{lemma}
\begin{proof} By Theorem 1 in \cite{Hironaka}, after dropping the contribution of higher order neighboring points.
\end{proof}
%%
%%

%%%%
%%%%
%%%%
\subsection{A canonical genus bound} 

The following result will be useful to control the genus of the components of certain canonical divisors on surfaces. 

\begin{lemma}[Canonical genus bound]
\source{chabauty-coleman-bound-surfaces}
\notready \label{LemmaCanonicalGenusBd} Let $S$ be a smooth projective irreducible surface over $k$. Let $D$ be a canonical divisor on $S$. Assume that $D>0$ and that $D$ is numerically effective. Let $C_1,...,C_\ell$ be the irreducible components of the support of $D$ and let $a_1,...,a_\ell$ be positive integers defined by $D=a_1C_1 +...+a_\ell C_\ell$. Then we have
\source{chabauty-coleman-bound-surfaces}
$$
\sum_{j=1}^\ell a_j\cdot (\gfrak_a(C_j)-1)\le c_1^2(S).
$$
\end{lemma}
%%
%%
\begin{proof} Let us rearrange the indices so that for certain $0\le r\le \ell$ we have $C_j^2>0$ for each $j\le r$ and $C_j^2\le 0$ for each $j>r$. By the adjunction formula we find
%%
\begin{equation} \label{EqnCanBd1} 
\source{chabauty-coleman-bound-surfaces}
\begin{aligned}
2\sum_{j=1}^\ell a_j\cdot (\gfrak_a(C_j)-1)&=\sum_{j=1}^\ell a_j\cdot ((D.C_j) +C_j^2) = D^2 + \sum_{j=1}^\ell a_j\cdot  C_j^2\\
&\le D^2 + \sum_{j=1}^r a_j \cdot C_j^2 \le D^2 + \sum_{j=1}^r a_j^2 \cdot C_j^2.
\end{aligned}
\end{equation}
%%
For any given $j$ we have $(D-a_jC_j).C_j=(\sum_{i\ne j} a_iC_i).C_j\ge 0$ since the $C_i$ are different irreducible curves. Also, $D.C_j\ge 0$ since $D$ is nef. Therefore
$$
\begin{aligned}
D^2=\sum_{j=1}^\ell a_j D.C_j = \sum_{j=1}^r a_j^2 \cdot C_j^2 + \sum_{j=1}^r  a_j\cdot (D-a_jC_j).C_j + \sum_{j=r+1}^\ell a_j D.C_j\ge \sum_{j=1}^r a_j^2 \cdot C_j^2 .
\end{aligned}
$$
Together with \eqref{EqnCanBd1}, this gives $2\sum_{j=1}^\ell a_j\cdot (\gfrak_a(C_j)-1)\le 2\cdot D^2=2c_1^2(X)$.
\end{proof}
%%
%%

%%%%%
%%%%%
%%%%%
%%%%%

\subsection{Branches} Let $S$ be a surface over $k$ and let $x\in S(k)$ be a smooth point of $S$. Given local parameters $s_1,s_2\in \mfrak_{S,x}$, the completed local ring at $x$ can be presented as $\widehat{\Ocal}_{S,x}\simeq  k[[s_1,s_2]]$.


 Let $C\subseteq S$ be a closed irreducible curve passing through $x$ with local equation $F\in \mfrak_{S,x}$. Let us factor $F=F_1\cdots F_r$ where each $F_j\in \widehat{\Ocal}_{S,x}$ is irreducible. 
%%
%% 
\begin{lemma}
\source{chabauty-coleman-bound-surfaces}
\notready For each $i\ne j$, the power series $F_i$ and $F_j$ are non-associated.
\end{lemma} 
\begin{proof} The special case $\Ocal_{S,x}\simeq k[s_1,s_2]_{(s_1,s_2)}$ is in \cite{KunzBranches}  Theorem 16.5.  As observed in \cite{GFthesis} Remark 2.58, the general case follows from a result of Chevalley. Indeed, $k$ is perfect since it is algebraically closed. By \cite{ZariskiSamuel2} VIII.13 Theorem 31, the local ring $\Ocal_{C,x}=\Ocal_{S,x}/(F)$ is analytically unramified. Thus, the completion $\widehat{\Ocal}_{C,x} = \widehat{\Ocal}_{S,x}/(F)$ does not have nilpotent elements.
\end{proof}
%%
%%
Hence, the prime ideals $\qfrak_j=(F_j)\subseteq \widehat{\Ocal}_{S,x}$ are different. They will be called \emph{branches} of $C$ at $x$.

Let $\nu: \widetilde{C}\to S$ be the composition of the normalizaton map $\widetilde{C}\to C$ and the inclusion $C\to S$. For each $y\in \nu^{-1}(x)$ we have a map induced on local rings $\nu^\#_y: \Ocal_{S,x}\to \Ocal_{\widetilde{C},y}$ which extends continuously to completed local rings $\widehat{\nu}^\#_y: \widehat{\Ocal}_{S,x}\to \widehat{\Ocal}_{\widetilde{C},y}$. We will need the following:

%%
%%
\begin{lemma}[Relation between normalization and branches]
\source{chabauty-coleman-bound-surfaces}
\notready \label{LemmaBranches} With the previous notation, there is a bijection between the set of branches of $C$ at $x$ and the points $y\in \nu^{-1}(x)$. More precisely, the ideals $\ker(\widehat{\nu}^\#_y)\subseteq \widehat{\Ocal}_{S,x}$ for $y\in \nu^{-1}(x)$ are exactly the branches of $C$ at $x$.
\source{chabauty-coleman-bound-surfaces}
\end{lemma}
%%
%%
\begin{proof} In the special case $\Ocal_{S,x}\simeq k[s_1,s_2]_{(s_1,s_2)}$ this is \cite{KunzBranches} Theorem 16.14; a similar argument works in general. Alternatively, the general case is \cite{GFthesis} Theorem 2.69 assuming that $k$ has characteristic $0$. This assumption is needed in other parts of \cite{GFthesis}, but not  for this particular result. In fact, the argument in \cite{GFthesis} reduces this result to \cite{DKatz} Corollary 5, where characteristic $0$ is not required.
\end{proof}
%%
%%

%%
%%
\begin{lemma}[Bound for fibres of normalization]
\source{chabauty-coleman-bound-surfaces}
\notready \label{LemmaBounddelta} With the previous notation, and assuming that $S$ is a projective surface, we have $\#\nu^{-1}(x)\le \delta(C,x)+1$.
\source{chabauty-coleman-bound-surfaces}
\end{lemma}
\begin{proof} From Lemma \ref{LemmaMultFormula}, it follows that the number of branches of $C$ at $x$ is at most $\mult_x(C)$. By Lemma \ref{LemmaBranches} we get $\nu^{-1}(x)\le \mult_x(C)$. The result follows from Lemma \ref{LemmaHironakaMult}.
\end{proof}
%%
%%


%%
%%
As the isomorphism $\widehat{\Ocal}_{S,x}\simeq  k[[s_1,s_2]]$ allows us to study branches by means of power series computations, the following result will be useful. It is a direct consequence of \cite{Ploski} Theorem 2.1.
%%
%%
\begin{lemma}
\source{chabauty-coleman-bound-surfaces}
\notready \label{LemmaPloski} Let $f(x_1,x_2)\in k[[x_1,x_2]]$ be an irreducible power series. Let $\phi_1(t),\phi_2(t)\in t\cdot k[[t]]$ be such that $f(\phi_1,\phi_2)=0$. Then $\ord_t\phi_1(t)\ge \ord_{x_2} f(0,x_2)$ and $\ord_t\phi_2(t)\ge \ord_{x_1} f(x_1,0)$.
\source{chabauty-coleman-bound-surfaces}
\end{lemma}
%%
%%



%%%%%
%%%%%
%%%%%
%%%%%

\subsection{Surfaces of general type} Let $S$ be a smooth, irreducible, projective surface over $k$.  The surface $S$ is of \emph{general type} if the canonical sheaf $\Omega^2_{S/k}$ is big, i.e. if $\dim_kH^0(S,(\Omega^2_{S/k})^{\otimes n})$ grows quadratically on $n$. More generally, if $X$ is a smooth, geometrically irreducible, projective surface over $L$, we say that $X$ is of general type if $X\otimes L^{alg}$ is a surface of general type over $L^{alg}$.

The following result is classical. Here, ``minimal'' means that $S$ contains no $(-1)$-curves (i.e. smooth rational curves with self-intersection $-1$).

%%
%%
\begin{lemma}
\source{chabauty-coleman-bound-surfaces}
\notready \label{LemmaNef} If $S$ is a minimal smooth projective surface of general type over $k$, then $c_1^2(S)\ge 1$ and the canonical sheaf $\Omega^2_{S/k}$ is nef. 
\source{chabauty-coleman-bound-surfaces}
\end{lemma}
%%
%%

The following simple condition ensures that the canonical sheaf of a general type surface is ample.

%%
%%
\begin{lemma}
\source{chabauty-coleman-bound-surfaces}
\notready\label{LemmaGTA} Let $S$ be a smooth projective surface of general type. If $S$ contains no smooth rational curves with self-intersection $-1$ or $-2$, then the canonical sheaf of $S$ is ample.  In particular, this is the case when $S$ contains no smooth curve of genus $0$.
\source{chabauty-coleman-bound-surfaces}
\end{lemma}
\begin{proof} This result is well-known. The argument is similar to that in \cite{MumfordApp}. We present the proof for the convenience of the reader.

 Let $K$ be a canonical divisor on $S$. Since $S$ does not contain $(-1)$-curves, it is minimal. Lemma \ref{LemmaNef} implies that $K^2>0$ and that $K$ is nef. Let $C\subseteq X$ be an irreducible curve. By the Nakai-Moishezon criterion, it suffices to show that $K.C\ne 0$. 

For the sake of contradiction, suppose that $K.C=0$. As $K^2>0$, the Hodge Index Theorem gives that either $C^2<0$ or $C$ is numerically trivial. The latter case cannot happen because $C.H>0$ for any ample divisor $H$, so we have $C^2<0$. By the adjunction formula $2\gfrak_a(C)-2 = C^2 +K\cdot C = C^2< 0$, which gives $\gfrak_a(C) = 0$. Thus $C$ is smooth and rational by Lemma \ref{LemmaHironakaGenus}. Furthermore, $C^2=C^2 +K\cdot C=2\gfrak_a(C)-2=-2$. Contradiction. So, $K$ is ample.
\end{proof}
%%
%%
Let us recall that a compact complex manifold $M$ is Brody hyperbolic if every holomorphic map $h:\C\to M$ is constant. As we are in the compact case, this is equivalent to Kobayashi hyperbolicity by Brody's theorem \cite{Brody}, and we simply say that $M$ is hyperbolic. It is conjectured that if $V$ is a smooth projective variety over $\C$ and  $V(\C)$ is hyperbolic, then $V$ is of general type. This is known for surfaces; see \cite{BHPV} Chapter VIII, Theorem 24.1.
%%
%%
\begin{lemma}
\source{chabauty-coleman-bound-surfaces}
\notready\label{LemmaGenTypeHyp} Let $S$ be a smooth, irreducible,  projective surface defined over $\C$. If $S(\C)$ is hyperbolic, then $S$ is of general type.
\source{chabauty-coleman-bound-surfaces}
\end{lemma}


%%
%%


\begin{lemma}
\source{chabauty-coleman-bound-surfaces}
\notready\label{LemmaGenTypeTranslate} Let $A$ be an abelian variety over $\C$ and let $X\subseteq A$ be a smooth projective surface. If $X$ is not an abelian surface and it does not contain elliptic curves, then $X$ is of general type.
\source{chabauty-coleman-bound-surfaces}
\end{lemma}
\begin{proof}
By \cite{Yau} Theorems 1 and 1'. Alternatively:  $X$ is hyperbolic \cite{Green}, then apply  Lemma \ref{LemmaGenTypeHyp}.
\end{proof}

%%%%%
%%%%%
%%%%%
%%%%%




\subsection{Point counting over finite fields} \label{SecZeta}
\source{chabauty-coleman-bound-surfaces}

In this paragraph we let $\kk$ be a finite field of characteristic $p$ with $q$ elements. Let $D$ be a projective curve defined over $\kk$; we do not require that $D$ be irreducible, connected, or smooth. Let $\kk'$ be a finite extension of $\kk$ such that all geometric irreducible components of $D$ are defined over $\kk'$, and all singular points of $D$ as well as all points of the normalization of $D\otimes \kk'$ mapping to them are $\kk'$-rational.

Let $C_1,...,C_r$ be the geometric irreducible components of $D$, which are defined over $\kk'$. Let $\nu_j:\widetilde{C}_j\to D\otimes \kk'$ be the normalization of $C_j$ composed with the inclusion $C_j\to D\otimes \kk'$. 

Our goal is to show the following estimate

%%
%%
\begin{lemma}
\source{chabauty-coleman-bound-surfaces}
\notready\label{LemmaWeilBound} With the previous notation, we have
\source{chabauty-coleman-bound-surfaces}
$$
\sum_{x\in D(\kk)} \sum_{j=1}^r \#\{y\in \widetilde{C}_j(\kk^{alg}) : \nu_j(y)=x\} \le (q+1)r + 2q^{1/2}\cdot \sum_{j=1}^r \gfrak_g(C_j).
$$
\end{lemma}
%%
%%
More generally, for a finite extension $L/\kk$ it will be convenient to define 
$$
A_D(L)=\sum_{x\in D(L)} \sum_{j=1}^r \#\{y\in \widetilde{C}_j(\kk^{alg}) : \nu_j(y)=x\}
$$
and $a_D(L)=A_D(L)-\#D(L)$. Note that $a_D(L)\ge 0$ and by the definition of $\kk'$ we see that the number $c_D:=a_D(\kk')$ satisfies $a_D(L)\le c_D$ for each finite extension $L/\kk$.

%%
\begin{lemma}
\source{chabauty-coleman-bound-surfaces}
\notready\label{LemmaLbig} If $L/\kk$ is a finite extension and $\kk'\subseteq L$, then $a_D(L)=c_D$ and
\source{chabauty-coleman-bound-surfaces}
$$
A_D(L)=\sum_{j=1}^r \#\widetilde{C}_j(L).
$$
\end{lemma}
%%
\begin{proof} The quantity $a_D(L)$ counts preimages of singularities of $D$ in the normalization, so it stabilizes when $L$ contains $\kk'$. Let $C$ be the disjoint union of the curves $C_j$ and let $\nu:C\to D\otimes \kk'$ be the morphism over $\kk'$ determined by all the maps $\nu_j$. Then $\nu$ is an isomorphism over $\kk'$ away from the preimages of singularities of $D$, all of which are $\kk'$-rational. Thus,  $A_D(L)=\#C(L)$.
\end{proof}


Note that $\#D(\kk)\le A_D(\kk) \le \# D(\kk) + c_D$. There is related work on the Weil conjectures for singular curves such as \cite{AubryPerret} and a natural approach to Lemma \ref{LemmaWeilBound} would be to use an upper bound on $\#D(\kk)$ together with an upper bound on $c_D$ coming from counting singularities. However, we will consider $\#D(L)+c_D$ at once, as this leads to sharper estimates.


For an integer $n\ge 1$ and $F$ a finite field of characteristic $p$, let $F_n\subseteq F^{alg}$ be the only extension of $F$ of degree $n$. Given a quasi-projective variety $V$ over $F$ (not necessarily smooth or irreducible) let $N_{V,F}(n)=\# V(F_n)$ and define the Zeta function on the variable $T$ by
$$
Z_{V,F}(T)=\exp\left(\sum_{n=1}^\infty \frac{N_{V,F}(n)}{n}\cdot T^n\right).
$$
In this generality, it is a theorem of Dwork \cite{Dwork} that $Z_{V,F}(T)\in \Q(T)$.

Let $N^*_{D,\kk}(n)=\# D(\kk_n)+c_D$ and observe that for all $n\ge 1$ we have 
%%
\begin{equation}\label{EqnBdAD}
\source{chabauty-coleman-bound-surfaces}
A_D(\kk_n)\le N^*_{D,\kk}(n). 
\end{equation}
%%
Let us define
$$
Z^*_{D,\kk}(T)=\exp\left(\sum_{n=1}^\infty \frac{N^*_{D,\kk}(n)}{n}\cdot T^n\right).
$$
Then $Z^*_{D,\kk}(T)=Z_{D,\kk}(T)/(1-T)^{c_D}\in \Q(T)$. Writing $d=[\kk':\kk]$, we have
%%
\begin{lemma}
\source{chabauty-coleman-bound-surfaces}
\notready\label{LemmaZprod} Let $\mu_d$ be the set of complex $d$-th roots of $1$. Then
\source{chabauty-coleman-bound-surfaces}
$$
\prod_{\zeta\in \mu_d} Z^*_{D,\kk}(\zeta\cdot T) = \prod_{j=1}^r Z_{\widetilde{C}_j,\kk'}(T^d).
$$
\end{lemma}
%%
\begin{proof} We note that for a given $n\ge 1$
$$
S_n:=\sum_{\zeta\in \mu_d}\frac{N^*_{D,\kk}(n)}{n}\cdot \zeta^nT^n=
\begin{cases}
0 & \mbox{ if } d\nmid n\\
d\cdot \frac{N^*_{D,\kk}(n)}{n}\cdot T^m & \mbox{ if } d|n.
\end{cases}
$$
In the second case we write $n=md$ and observe that $S_{md} = m^{-1}N^*_{D,\kk}(md) T^{md}$. Using $\kk_{md}=\kk'_m\supseteq \kk'$ and Lemma \ref{LemmaLbig} we get
$$
N^*_{D,\kk}(md) =\# D(\kk'_m) + c_D=\# D(\kk'_m)+ a_D(\kk'_m)=A_D(\kk'_m)=\sum_{j=1}^r N_{\widetilde{C}_j,\kk'}(m)
$$
which proves the result.
\end{proof}
%%

\begin{proof}[Proof of Lemma \ref{LemmaWeilBound}] The Weil conjectures for curves (proved by Weil) applied to the curves $\widetilde{C}_j$ over $\kk'$ together with Lemma \ref{LemmaZprod} show that $Z^*_{D,\kk}(T)$ on $\C$ has $r$ poles of modulus $q$, a total of $2\cdot \sum_{j=1}^r \gfrak_g(C_j)$ zeros of modulus $q^{1/2}$, $r$ poles of modulus $1$, and no other zeros or poles in $\C$. A standard computation taking the logarithmic derivative of $Z^*_{D,\kk}(T)$, comparing coefficients of $T^n$, and applying the triangle inequality, shows that for each $n\ge 1$ we have $N^*_{D,\kk}(n)\le r(q^n+1) + 2q^{n/2} \sum_{j=1}^r \gfrak_g(C_j)$. We conclude by \eqref{EqnBdAD} and taking $n=1$. 
\end{proof}







%%%%%%%%%%%%%%%%%%%%%%%%%%%%%%%%%%%%%%
%%%%%%%%%%%%%%%%%%%%%%%%%%%%%%%%%%%%%%
%%%%%%%%%%%%%%%%%%%%%%%%%%%%%%%%%%%%%%
%%%%%%%%%%%%%%%%%%%%%%%%%%%%%%%%%%%%%%
%%%%%%%%%%%%%%%%%%%%%%%%%%%%%%%%%%%%%%
%%%%%%%%%%%%%%%%%%%%%%%%%%%%%%%%%%%%%%
